\documentclass[11pt]{article}

\usepackage[utf8]{inputenc}
\usepackage[T1]{fontenc}
\usepackage[spanish]{babel}
\usepackage{lmodern}
\usepackage[a4paper,margin=2.5cm]{geometry}
\usepackage{amsmath}
\usepackage{booktabs}
\usepackage{graphicx}
\usepackage{pgfplots}
\pgfplotsset{compat=1.18}
\usepackage{hyperref}
\hypersetup{colorlinks=true, linkcolor=black, urlcolor=blue, citecolor=black}

\title{Base de Datos II (CS2042)\\
Proyecto del Curso: Gestor de Bases de Datos Multimodal\\
Entregable 1}
% Editar aquí si hay más miembros en el equipo:
\author{Brando Lopez}
\date{Universidad de Ingeniería y Tecnología (UTEC)\\
Docente: Percy Lovon Ramos\\
Ciclo 2026-II}

\begin{document}

\maketitle
\tableofcontents
\newpage

\section{Introducción y objetivos}

MiniDB es un motor de bases de datos construido \emph{desde cero} en Python,
sin librerías de bases de datos ni ORMs, como exige el proyecto del curso.
Toda la persistencia se implementa directamente sobre archivos binarios
organizados en \textbf{páginas de 4 KB} leídas y escritas con
\texttt{seek}/\texttt{read}/\texttt{write} y serialización \texttt{struct}.

El objetivo central es comprender y demostrar experimentalmente cómo la
\emph{organización física} de los datos determina el costo de las operaciones,
medido en términos de lo que realmente importa en un DBMS: el número de
\textbf{bloques de disco} leídos y escritos (I/O), no solo el tiempo de CPU.

Los objetivos específicos del Entregable 1 son:

\begin{itemize}
  \item Implementar un heap file sobre páginas ranuradas con lista de libres.
  \item Implementar un \emph{Sequential File} (área principal ordenada + área
        de overflow encadenada + reorganización con fill factor).
  \item Implementar índices B+ Tree y Hash Extensible sobre archivo,
        y un R-Tree como extensión espacial.
  \item Integrar un motor SQL mínimo (parser, catálogo, ejecutor con plan)
        que elija la vía de acceso adecuada.
  \item Instrumentar un contador de I/O físico y validar experimentalmente
        las complejidades teóricas: $O(P)$ del escaneo secuencial,
        $O(\log N)$ del B+ Tree, $O(1)$ del hash, búsqueda binaria del
        archivo secuencial.
\end{itemize}

\section{Diseño físico de almacenamiento}

\subsection{Página ranurada (slotted page) de 4 KB}

La unidad básica de almacenamiento es la página de 4096 bytes con el
siguiente layout:

\begin{center}
\begin{tabular}{ll}
\toprule
\textbf{Región} & \textbf{Contenido} \\
\midrule
Cabecera (16 B) & \texttt{page\_id} (I), \texttt{slot\_count} (H),
  \texttt{record\_count} (H), \texttt{free\_space\_offset} (H),
  \texttt{next\_page\_id} (i), \texttt{prev\_page\_id} (i) \\
Slot array & 6 B por slot: (\texttt{offset} H, \texttt{length} H, \texttt{alive} B, res B) \\
\ldots & espacio libre contiguo \\
Tuplas & escritas desde el final de la página hacia atrás \\
\bottomrule
\end{tabular}
\end{center}

Los registros crecen desde el final de la página hacia el inicio y el slot
array crece desde el inicio; \texttt{free\_space\_offset} marca el límite
del área de datos. Los punteros \texttt{next\_page\_id}/\texttt{prev\_page\_id}
($-1$ si no hay vecino) encadenan las páginas de datos de un archivo como
una lista doblemente enlazada (en el overflow del Sequential File quedan
en $-1$: ahí el encadenamiento es a nivel de registro). La eliminación es
\textbf{lógica}: marca el slot como muerto
(\texttt{alive = 0}) sin mover bytes, de modo que el \texttt{slot\_id}
permanece estable. Si una inserción no encuentra espacio contiguo, la página
se compacta (reubicando solo los registros vivos) antes de declararla llena.

\subsection{RID}

Todo registro se identifica por \texttt{RID = (page\_id, slot\_id)}. El RID
es estable ante compactaciones y es la referencia que devuelven todos los
índices. En el Sequential File, a los RIDs del área de overflow se les suma
una base ($2^{20}$) para distinguir el área sin estructuras adicionales.

\subsection{Serialización con \texttt{struct}}

Las filas se serializan campo a campo según el esquema, sin pickle:

\begin{center}
\begin{tabular}{lll}
\toprule
\textbf{Tipo SQL} & \textbf{Formato struct} & \textbf{Tamaño} \\
\midrule
INT & \texttt{<i} & 4 B \\
FLOAT & \texttt{<d} & 8 B \\
BOOL & \texttt{<?} & 1 B \\
VARCHAR($n$) & \texttt{<$n$s} & $n$ B (padding NUL) \\
POINT & \texttt{<dd} & 16 B \\
TEXT & \texttt{<H} + bytes & variable \\
\bottomrule
\end{tabular}
\end{center}

Las claves de índice usan la misma codificación de longitud fija
(\texttt{encode\_key}/\texttt{decode\_key}), lo que permite comparar
valores decodificados y calcular capacidades exactas por página.

\subsection{Heap file con free list}

El heap file guarda una página 0 de cabecera (magic, contadores y la
\emph{free list} de RIDs de slots eliminados, con capacidad
\texttt{MAX\_FREE} $= (4096 - \text{cabecera}) / 6$ entradas). La inserción
reutiliza primero slots de la lista de libres; si está vacía intenta la
última página y, en última instancia, agrega una página nueva. Los slots
muertos que no caben en la lista se contabilizan (\texttt{dropped\_free})
y se recuperan con un escaneo de respaldo solo cuando hubo desborde, evitando
barridos $O(P)$ innecesarios.

\section{Sequential File}

\subsection{Estructura}

Dos archivos por tabla: \texttt{<tabla>.seq} (área principal, registros
ordenados físicamente por la clave primaria, con búsqueda binaria por
páginas) y \texttt{<tabla>.ovf} (área de overflow). Formato de registro en
ambas áreas:

\begin{center}
\texttt{[key\_bytes (fijo)][next\_page: I][next\_slot: H][payload]}
\end{center}

El puntero \texttt{next} mantiene el \emph{orden lógico} encadenando los
registros del overflow entre los de la principal. La cabecera del
\texttt{.seq} guarda además los punteros \texttt{head} y \texttt{tail} de la
cadena: \texttt{head} permite arrancar el recorrido tras eliminar el primer
registro, y \texttt{tail} da un camino rápido $O(1)$ para inserciones al
final (clave mayor que todas), decisión clave para que la carga masiva
ordenada sea $O(N)$ en vez de $O(N^2)$.

\subsection{Búsqueda}

La búsqueda puntual hace \textbf{búsqueda binaria por páginas} en la
principal (comparando la primera clave física de cada página, que sigue
siendo válida aunque el registro esté eliminado), retrocede hasta el último
registro \emph{vivo} con clave $\le$ la buscada y desde ahí recorre la
cadena del overflow, deteniéndose al pasar la clave. Los rangos
(\texttt{BETWEEN}, $<$, $\le$, $>$, $\ge$) comparten el mismo punto de
partida y recorren la cadena hasta salir del rango. Costo:
$O(\log P) + $ longitud de la cadena atravesada.

\subsection{Inserción y eliminación}

Insertar ubica al predecesor con la búsqueda anterior, escribe el registro
en el overflow (reutilizando una free list simple en la cabecera del
\texttt{.ovf}) y reenlaza: \texttt{new.next = pred.next; pred.next = new}.
Eliminar reenlaza al predecesor y marca el slot como muerto; el espacio del
overflow se reutiliza vía la free list y el de la principal se recupera en
la reorganización.

\subsection{Reorganización con fill factor}

\texttt{reorganize(fill\_factor)} recorre la cadena completa (vivos de
ambas áreas en orden), reescribe un \texttt{.seq} nuevo con páginas llenas
al 70--80\% (empaquetado greedy), reconstituye \texttt{next} como el
siguiente físico y trunca el \texttt{.ovf} a su cabecera. El reemplazo es
atómico (archivo temporal + \texttt{os.replace}) y la operación devuelve
estadísticas: páginas antes/después, registros movidos y muertos purgados.

\section{Índices}

\subsection{B+ Tree}

Persistido en su propio archivo de páginas de 4 KB: página 0 de cabecera
(magic, raíz, número de páginas, tamaño de clave) y una página por nodo,
con capacidades calculadas para que cada nodo quepa exactamente en una
página. Soporta claves duplicadas ordenando por la tupla \texttt{(key, rid)}
(los separadores internos también son pares). Las hojas forman una lista
doblemente enlazada con punteros \texttt{next\_leaf\_id} y
\texttt{prev\_leaf\_id}, de modo que un \emph{range search} baja una vez por
el árbol ($O(\log N)$) y luego recorre hojas contiguas en cualquier
sentido. Incluye caché de
nodos y modo de carga masiva con volcado diferido.

\subsection{Hash Extensible}

Función hash FNV-1a de 32 bits implementada a mano. Directorio de $2^g$
entradas u32 indexadas por los $g$ bits menos significativos del hash;
buckets de capacidad fija con profundidad local; split con duplicación del
directorio cuando la profundidad local iguala a la global.

\paragraph{Lección de ingeniería (EXH1 $\to$ EXH2).}
La primera versión del índice cargaba \emph{todo} el archivo a memoria al
abrirlo y lo reescribía completo en cada mutación. Como el ejecutor abre el
índice en cada consulta, una búsqueda puntual sobre 100 000 claves costaba
\textbf{259 lecturas de bloque y 60 ms}: el costo de lectura era $O(B)$
(páginas del índice) por consulta, no $O(1)$. La versión actual (EXH2)
reserva el directorio con capacidad fija desde el inicio (región
\emph{sparse} vía \texttt{truncate}, sin costo físico hasta usarse), ubica
cada bucket en un offset fijo (\texttt{bucket\_id} = posición física) y lee
perezosamente: cabecera + 1 página de directorio + 1 página de bucket.
Resultado medido: \textbf{5 lecturas y 0.20 ms} por búsqueda (sección 7.3).
El formato viejo no se soporta; los índices hash existentes se recrean.

\subsection{R-Tree (extensión espacial)}

Como adelanto de la Parte 2 del proyecto, el motor incluye un R-Tree sobre
páginas de 4 KB con split cuadrático de Guttman, búsqueda por radio con
poda por MBR y KNN con cola de prioridad por \emph{mindist}. Sirve como
índice secundario de columnas \texttt{POINT} y se integra al planificador
condicional (condiciones \texttt{IN (punto, radio)} y
\texttt{KNN (punto, k)}).

\section{Motor SQL y planificador}

\subsection{Subconjunto SQL}

Parser recursivo descendente propio (keywords case-insensitive):

\begin{itemize}
  \item \texttt{CREATE TABLE} (con PK de una columna y
        \texttt{USING HEAP|SEQUENTIAL}), \texttt{CREATE TABLE ... FROM FILE},
        \texttt{LOAD INTO ... FROM FILE}, \texttt{DROP TABLE}.
  \item \texttt{CREATE INDEX [nombre] ON t (col) USING BTREE|HASH|RTREE}.
  \item \texttt{INSERT INTO ... VALUES}, \texttt{DELETE FROM ... WHERE col = v}.
  \item \texttt{SELECT} con \texttt{WHERE} (=, $<$, $\le$, $>$, $\ge$,
        \texttt{BETWEEN}, condiciones espaciales) con conjunciones
        \texttt{AND} (el planificador elige el mejor índice de la
        conjunción y aplica el resto como filtro residual),
        \texttt{LIMIT/OFFSET} y
        agregados \texttt{COUNT/MIN/MAX/SUM/AVG}.
\end{itemize}

\subsection{Selección de la vía de acceso}

El ejecutor construye un plan con tiempos por paso y elige:

\begin{itemize}
  \item \textbf{Index Scan} (B+ Tree o Hash) para igualdad sobre columna
        indexada; \textbf{Index Range Scan} (B+ Tree) para rangos.
  \item \textbf{Binary Search (Sequential)} para predicados sobre la PK de
        una tabla \texttt{USING SEQUENTIAL} (el propio archivo ordenado es
        la vía de acceso; no se genera B+ Tree redundante).
  \item \textbf{Sequential Scan} cuando no hay índice usable.
\end{itemize}

\subsection{DiskCounter: medición honesta de I/O}

Toda la I/O física del motor pasa por objetos archivo envueltos en
\texttt{CountedFile}, que suma 1 lectura por \texttt{read()} no vacío y 1
escritura por \texttt{write()} con datos (1 llamada = 1 bloque, pues el
código siempre lee/escribe páginas completas de 4 KB). Cada consulta crea un
contador fresco y devuelve \texttt{result["io"] = \{reads, writes\}}.

La métrica es honesta por construcción: cuenta \emph{llamadas físicas
reales} al archivo; las lecturas servidas desde las cachés de páginas en
memoria no cuentan, exactamente como un buffer pool evita I/O real. Así, los
números reportados son comparables con el modelo de costos del curso.

\section{Arquitectura del sistema}

\begin{itemize}
  \item \textbf{Motor} (Python puro): \texttt{storage/} (página, heap,
        sequential file, serialización, disk counter), \texttt{indexes/}
        (B+ Tree, Hash Extensible, R-Tree), \texttt{engine/} (parser,
        catálogo JSON, ejecutor, cargador CSV).
  \item \textbf{API REST} (FastAPI): \texttt{/api/query},
        \texttt{/api/tables}, \texttt{/api/tables/\{t\}/reorganize}, carga e
        inferencia de CSV.
  \item \textbf{Frontend} (React + Vite + Leaflet): consola SQL con plan e
        I/O por consulta, explorador de tablas y mapa espacial.
  \item \textbf{Despliegue}: \texttt{docker compose} (backend + frontend con
        nginx), volumen persistente para los datos.
\end{itemize}

\section{Experimentación}

\subsection{Metodología}

Scripts en \texttt{benchmarks/} que usan el motor directamente (sin HTTP)
sobre datos sintéticos \texttt{(id INT PK, nombre VARCHAR(30), valor FLOAT)}
con semillas aleatorias fijas para reproducibilidad. Se mide por consulta la
I/O física (\texttt{io.reads}/\texttt{io.writes} del DiskCounter) y la
latencia en ms; se reporta promedio $\pm$ desviación estándar. La carga es
en lote (\texttt{bulk\_load\_rows}, el mismo camino físico de la carga CSV),
no tupla a tupla. Las variantes con índice barajan las claves; el Sequential
File se carga ordenado (camino rápido del \texttt{tail}) y se reorganiza
antes de los experimentos de búsqueda, que es la operación periódica normal
de esta organización. Resultados crudos en \texttt{benchmarks/results/}.

\subsection{Experimento 1: inserción masiva}

Inserción de $N$ registros por organización; tiempo total, escrituras de
bloques de 4 KB y filas/segundo.

\begin{center}
\begin{tabular}{lrrrr}
\toprule
\textbf{Organización} & $N{=}10^3$ & $N{=}10^4$ & $N{=}5\cdot10^4$ & $N{=}10^5$ \\
\midrule
Heap sin índice      & 0.004 & 0.045 & 0.228 & 0.476 \\
Heap + B+ Tree (PK)  & 0.019 & 0.234 & 1.404 & 3.590 \\
Heap + Hash Ext.\ (PK) & 0.016 & 0.195 & 1.094 & 2.545 \\
Sequential File      & 0.006 & 0.066 & 0.334 & 0.674 \\
\bottomrule
\end{tabular}

\smallskip
(Tiempos en segundos.)
\end{center}

\begin{center}
\begin{tikzpicture}
\begin{loglogaxis}[
    width=0.85\textwidth, height=6.5cm,
    xlabel={$N$ (registros)}, ylabel={Tiempo total (s)},
    legend pos=north west, grid=major, mark size=1.8pt]
\addplot coordinates {(1000,0.004) (10000,0.045) (50000,0.228) (100000,0.476)};
\addplot coordinates {(1000,0.019) (10000,0.234) (50000,1.404) (100000,3.590)};
\addplot coordinates {(1000,0.016) (10000,0.195) (50000,1.094) (100000,2.545)};
\addplot coordinates {(1000,0.006) (10000,0.066) (50000,0.334) (100000,0.674)};
\legend{Heap sin índice, Heap + B+ Tree, Heap + Hash Ext., Sequential File}
\end{loglogaxis}
\end{tikzpicture}
\end{center}

Para $N = 100\,000$, las escrituras de bloques fueron: heap 1192, B+ Tree
1543, Hash 1450 y Sequential 1432 (210 168, 27 857, 39 287 y 148 331
filas/s respectivamente).

\paragraph{Análisis.}
El heap pelado es el límite inferior teórico: escribe cada página de datos
una sola vez (1192 escrituras $\approx$ $N \cdot 48$ B $/ 4096$ B) y su
throughput es constante ($\sim$220k filas/s en todo el rango), es decir
$O(1)$ amortizado por tupla. Los índices pagan el mantenimiento de la
estructura: el B+ Tree escribe $\sim$30\% más bloques y es $7.5\times$ más
lento en tiempo a $N{=}10^5$ (splits y propagación, $O(\log N)$ por tupla);
el hash queda intermedio. El Sequential File con carga ordenada es el
segundo más rápido (0.67 s): gracias al puntero \texttt{tail} persistido,
cada inserción es $O(1)$ y solo escribe páginas del área de overflow; sin
ese camino rápido la inserción ordenada recorrería la cadena y degradaría a
$O(N^2)$.

\subsection{Experimento 2: búsquedas puntuales}

Búsquedas de igualdad \texttt{WHERE id = k} sobre $N = 100\,000$ registros,
con claves aleatorias existentes (1000 consultas; 100 para el full scan por
su costo de $\sim$0.2 s por consulta).

\begin{center}
\begin{tabular}{lrrr}
\toprule
\textbf{Organización} & \textbf{Reads prom.} & \textbf{Latencia prom.\ (ms)} \\
\midrule
Heap sin índice (full scan) & $1192 \pm 0$    & $183.5 \pm 9.4$ \\
Heap + B+ Tree (PK)         & $6.01 \pm 0.09$ & $0.26 \pm 0.03$ \\
Heap + Hash Ext.\ (PK)      & $5 \pm 0$       & $0.20 \pm 0.02$ \\
Sequential File (binaria)   & $16.94 \pm 0.26$ & $0.24 \pm 0.02$ \\
\bottomrule
\end{tabular}
\end{center}

\paragraph{Análisis.}
Los números replican la teoría casi exactamente. El full scan lee las 1192
páginas de datos en \emph{todas} las consultas ($O(P)$, std = 0: no depende
de la clave). El B+ Tree lee $O(\log N)$: cabecera + camino raíz$\to$hoja
(2 niveles bastan para 100k claves de 4 B) + cabecera del heap + 1 página de
datos = 6 bloques. El hash gana la igualdad con $O(1)$: cabecera + 1 página
de directorio + 1 página de bucket + 2 del heap = 5 bloques y la menor
latencia (0.20 ms) — el acceso es directo por función hash, sin
comparaciones en niveles intermedios. La búsqueda binaria del Sequential
File cuesta $\sim$17 lecturas ($\log_2 P$ páginas de la principal más la
caminata corta y el fetch), ligeramente por encima del B+ Tree porque la
binaria compara una clave por página sin fan-out interno; sigue siendo dos
órdenes de magnitud mejor que el scan. Vale contrastar con la versión
inicial del hash (EXH1): 259 lecturas y 60 ms por la misma búsqueda — el
formato de archivo importa tanto como la estructura lógica.

\subsection{Experimento 3: búsquedas por rango con selectividad variable}

Rangos \texttt{WHERE id BETWEEN a AND b} del 0.1\%, 1\%, 5\%, 10\% y 25\%
del total, sobre $N = 100\,000$ (50 consultas por celda; 10 en el full
scan). Las consultas llevan \texttt{LIMIT} explícito para medir el barrido
completo del rango sin el tope de seguridad de la consola.

\begin{center}
\small
\begin{tabular}{lrrrr}
\toprule
\textbf{Organización} & \textbf{Sel.\ \%} & \textbf{Tuplas} & \textbf{Reads prom.} & \textbf{ms prom.} \\
\midrule
Heap sin índice & 0.1--25 & 100--25 000 & $1192 \pm 0$ & 178--212 \\
\midrule
Heap + B+ Tree & 0.1 & 100    & $102 \pm 1.7$   & $3.2 \pm 3.5$ \\
Heap + B+ Tree & 1   & 1 000  & $686 \pm 10.3$  & $22.2 \pm 8.8$ \\
Heap + B+ Tree & 5   & 5 000  & $1197 \pm 4.8$  & $50.9 \pm 12.3$ \\
Heap + B+ Tree & 10  & 10 000 & $1229 \pm 4.1$  & $66.3 \pm 13.5$ \\
Heap + B+ Tree & 25  & 25 000 & $1278 \pm 6.0$  & $111.7 \pm 11.9$ \\
\midrule
Sequential File & 0.1 & 100    & $18 \pm 0.6$   & $0.52 \pm 0.04$ \\
Sequential File & 1   & 1 000  & $51 \pm 1.2$   & $3.3 \pm 2.6$ \\
Sequential File & 5   & 5 000  & $204 \pm 1.5$  & $17.0 \pm 6.4$ \\
Sequential File & 10  & 10 000 & $396 \pm 1.5$  & $36.0 \pm 9.4$ \\
Sequential File & 25  & 25 000 & $972 \pm 1.5$  & $103.8 \pm 14.6$ \\
\bottomrule
\end{tabular}
\end{center}

\begin{center}
\begin{tikzpicture}
\begin{semilogxaxis}[
    width=0.85\textwidth, height=6.5cm,
    xlabel={Selectividad del rango (\% del total)},
    ylabel={Lecturas de bloques (prom.)},
    legend pos=north west, grid=major, mark size=1.8pt]
\addplot coordinates {(0.1,1192) (1,1192) (5,1192) (10,1192) (25,1192)};
\addplot coordinates {(0.1,101.56) (1,686.28) (5,1196.64) (10,1228.52) (25,1278)};
\addplot coordinates {(0.1,18.24) (1,51.42) (5,204.22) (10,396.02) (25,972.22)};
\legend{Heap sin índice (scan), Heap + B+ Tree, Sequential File}
\end{semilogxaxis}
\end{tikzpicture}
\end{center}

\paragraph{Análisis: el punto de cruce.}
El full scan es insensible a la selectividad (siempre 1192 páginas: hay que
barrer todo el archivo). El B+ Tree degrada con el tamaño del rango porque
los RIDs que devuelve están \emph{dispersos} por todo el heap (la carga fue
desordenada): cada tupla del rango puede tocar una página de datos distinta.
El \textbf{punto de cruce en I/O está entre 1\% y 5\%}: a 1\% el índice lee
686 páginas y gana claramente; a 5\% lee 1197, ya \emph{más} que barrer el
archivo completo. En latencia el scan nunca gana en este experimento (112 ms
del B+ Tree vs 212 ms del scan a 25\%) porque el scan además deserializa las
100 000 tuplas: el CPU de deserialización domina sobre la I/O extra — un
matiz útil: el cruce en bloques no implica cruce en tiempo cuando el plan
alternativo paga CPU por tupla. El Sequential File domina los rangos en toda
selectividad medida: al ser una organización \emph{clustered} el rango es
físicamente contiguo, así que sus lecturas escalan con la fracción del
archivo ($972 \approx 82\%$ de las páginas a 25\%, por el fill factor del
75\%) sin saltos; a 0.1\% es $5.6\times$ más económico que el B+ Tree (18 vs
102 lecturas). Es la confirmación experimental de por qué los rangos sobre
la clave de clustering son el caso ideal de las organizaciones ordenadas.

\subsection{Experimento 4: sensibilidad al tamaño de bloque}

B+ Tree aislado (sin heap) con $N = 100\,000$ claves \texttt{INT} de 4
bytes insertadas en orden aleatorio, variando el tamaño de página
$B \in \{1024, 2048, 4096, 8192\}$ bytes (parámetro \texttt{page\_size}
del \texttt{BPlusTree}; el resto del motor sigue en 4 KB). Para cada $B$
se mide la capacidad máxima de cada tipo de nodo según su formato de
serialización, la altura resultante, las lecturas de bloques de 1000
búsquedas puntuales con la caché de nodos vaciada antes de cada consulta
(peor caso: cada nivel del camino raíz$\to$hoja es una lectura física) y
el tamaño total del archivo del índice.

\begin{center}
\small
\begin{tabular}{rrrrrrr}
\toprule
\textbf{$B$ (B)} & \textbf{Cap.\ hoja} & \textbf{Cap.\ interno} &
\textbf{Altura} & \textbf{Páginas} & \textbf{MiB} & \textbf{Reads prom.} \\
\midrule
1024 & 101 & 72  & 3 & 1449 & 1.42 & $3.03 \pm 0.16$ \\
2048 & 203 & 145 & 3 & 708  & 1.38 & $3.01 \pm 0.11$ \\
4096 & 408 & 292 & 3 & 351  & 1.37 & $3.01 \pm 0.08$ \\
8192 & 818 & 584 & 2 & 160  & 1.25 & $2.00 \pm 0.00$ \\
\bottomrule
\end{tabular}
\end{center}

\begin{center}
\begin{minipage}{0.48\textwidth}
\centering
\begin{tikzpicture}
\begin{loglogaxis}[
    width=\textwidth, height=6cm,
    log basis x={2}, log basis y={2},
    xlabel={$B$ (bytes)}, ylabel={Capacidad (entradas/nodo)},
    xtick={1024,2048,4096,8192}, xticklabels={1K,2K,4K,8K},
    legend pos=north west, grid=major, mark size=1.8pt]
\addplot coordinates {(1024,72) (2048,145) (4096,292) (8192,584)};
\addplot coordinates {(1024,101) (2048,203) (4096,408) (8192,818)};
\legend{Nodo interno, Nodo hoja}
\end{loglogaxis}
\end{tikzpicture}
\end{minipage}\hfill
\begin{minipage}{0.48\textwidth}
\centering
\begin{tikzpicture}
\begin{semilogxaxis}[
    width=\textwidth, height=6cm,
    log basis x={2},
    xlabel={$B$ (bytes)}, ylabel={Altura $h$},
    xtick={1024,2048,4096,8192}, xticklabels={1K,2K,4K,8K},
    ytick={2,3}, ymin=1.5, ymax=3.5,
    grid=major, mark size=1.8pt]
\addplot coordinates {(1024,3) (2048,3) (4096,3) (8192,2)};
\end{semilogxaxis}
\end{tikzpicture}
\end{minipage}
\end{center}

\paragraph{Análisis.}
Duplicar $B$ duplica casi exactamente el fan-out: la capacidad del nodo
interno crece $72 \to 145 \to 292 \to 584$ (separadores de 10 B más un
hijo por página) y la de la hoja $101 \to 203 \to 408 \to 818$ (entradas
de 10 B) — el overhead fijo de cabecera (11 B) se vuelve despreciable
desde 2 KB. Con inserción aleatoria las hojas quedan al $\sim$70\% de
ocupación, así que $h = \lceil \log_{f}(N/(\text{hoja} \cdot 0.7))
\rceil$ baja de 3 a 2 recién en $B = 8$ KB: con $f \approx 584$ un solo
nivel interno basta para direccionar las $\sim$175 hojas. Las lecturas de
búsqueda replican la altura página a página ($3.0$ vs $h = 3$; el $0.03$
extra corresponde a las consultas cuya clave es la máxima de su hoja y
leen la hoja siguiente para confirmar que no hay duplicados), confirmando
que el costo de la igualdad es exactamente $\Theta(h)$ y que achatar el
árbol con páginas grandes se traduce en I/O real ahorrada. El tamaño del
índice es prácticamente constante ($1.25$--$1.42$ MiB): las páginas
grandes amortiguan cabeceras (160 páginas vs 1449 para la misma
información). El costo aparece del otro lado: la carga tarda $3.7$ s con
8 KB frente a $1.6$ s con 1 KB (serializar y escribir nodos de 8 KB por
split es más caro) y, en un sistema real, leer un bloque de 8 KB para
usar una entrada de 14 B desperdicia ancho de banda si la carga es de
búsquedas puntuales. La elección de $B$ es entonces un compromiso entre
altura del índice y granularidad de transferencia; 4 KB coincide con la
página del sistema operativo y equilibra ambos extremos.

\section{Conclusiones}

\begin{itemize}
  \item Las complejidades teóricas se verificaron experimentalmente con I/O
        física real: $O(P)$ del scan (1192 lecturas constantes), $O(\log N)$
        del B+ Tree (6 lecturas), $O(1)$ del hash extensible (5 lecturas) y
        búsqueda binaria del archivo secuencial ($\sim$17 lecturas).
  \item El formato de persistencia importa tanto como la estructura de
        datos: reescribir el índice hash completo por operación lo hacía
        $\sim$$300\times$ más lento en lectura que su versión paginada perezosa.
  \item No hay estructura universalmente mejor: el heap gana en inserción,
        el hash en igualdad, el B+ Tree equilibra igualdad y rangos
        moderados, y el archivo secuencial domina los rangos gracias a la
        contigüedad física — el planificador existe precisamente para elegir.
  \item La selectividad determina cuándo un índice deja de convenir: el
        punto de cruce medido entre el Index Range Scan y el full scan está
        entre 1\% y 5\% para esta relación, coherente con lo que estiman los
        optimizadores de costos de los DBMS reales.
  \item El tamaño de bloque controla el fan-out y por tanto la altura del
        B+ Tree: duplicar la página duplica las entradas por nodo, y con 8 KB
        el árbol de 100\,000 claves baja de 3 a 2 niveles, reduciendo la
        búsqueda puntual de 3 a 2 lecturas físicas; el índice ocupa el mismo
        espacio, pero la construcción se encarece.
\end{itemize}

\end{document}
